% !TEX root = hazel2020sensors.tex
% FICTIONAL TEST MATERIAL: this paper does not exist; it was written for the
% citation-audit fixture.
\documentclass{article}
\usepackage{amsmath}

\title{Low-Cost Pressure Sensors for Transient Monitoring}
\author{L. Fixture \and M. Fixture}
\date{Journal of Fixture Studies, 2020}

\begin{document}
\maketitle

\noindent\textbf{FICTIONAL TEST MATERIAL: this paper does not exist; it was
written for the citation-audit fixture.} No author, utility, venue, or
measurement described here is real.

\begin{abstract}
Pressure transients in water mains last a fraction of a second, which is far
shorter than the sampling interval of a typical network pressure logger. We
built low-cost pressure sensors that sample at 100~Hz and installed them beside
standard 1~Hz loggers at eight sites of a fictional network. The 100~Hz sensors
captured pressure transients above 5~m head that the 1~Hz loggers missed. Each
sensor cost a small fraction of a commercial transient logger. We describe the
hardware, the side-by-side trial, and the transient events recorded.
\end{abstract}

\section{Introduction}
\label{sec:introduction}

A pressure transient is a short pressure wave caused by a sudden change in
flow, such as a pump trip or a fast valve closure. Large or repeated transients
fatigue pipes and joints. Standard network loggers record pressure once per
second or less often, so a transient that rises and falls within a second can
pass between two samples and leave no trace.

Commercial transient loggers sample fast enough but are too expensive to deploy
widely. This paper asks whether a low-cost sensor sampling at 100~Hz can
capture the transients that the standard loggers miss.

\section{Method}
\label{sec:method}

\subsection{Sensor design}
\label{sec:method-sensor}

Each sensor combines an industrial pressure transducer, a small
microcontroller, and a battery. The microcontroller samples the transducer at
100~Hz and stores any one-second window in which the pressure changes by more
than 2~m head. The parts for one sensor cost less than one tenth of the price
of a commercial transient logger.

\subsection{Side-by-side trial}
\label{sec:method-trial}

At eight sites of the fictional Fixture Meadows network, a 100~Hz sensor was
installed on the same hydrant tapping as an existing 1~Hz logger. Both devices
ran for twelve weeks. Their clocks were synchronised, so each event recorded by
one device could be checked against the record of the other.

\section{Results}
\label{sec:results}

\subsection{Transients captured}
\label{sec:results-captured}

The low-cost sensors sampling at 100~Hz captured pressure transients above 5~m
head that the 1~Hz loggers missed. Over the twelve weeks, the 100~Hz sensors
recorded 37 transients larger than 5~m head across the eight sites. The 1~Hz
loggers showed no excursion above 5~m head for 31 of these 37 events; for the
other six, the 1~Hz record showed only a small step.

\subsection{Sources of the transients}
\label{sec:results-sources}

Most of the transients coincided with pump starts and stops at a booster
station. The others matched fast hydrant closures during flushing work.

\subsection{Cost and reliability}
\label{sec:results-cost}

All eight sensors ran for the full twelve weeks on one battery. At their low
cost, a utility could deploy them at many more sites than it could with
commercial transient loggers.

\section{Limitations}
\label{sec:limitations}

This study tests transient capture only.

This study does not examine bursts: we did not relate the recorded transients
to main breaks, and no burst occurred at the trial sites during the twelve
weeks. This study does not examine metering: the sensors measure pressure, not
flow, and no flow or consumption data were used. The trial covered eight sites
for twelve weeks, so the number of transients is small.

\end{document}
